\documentclass[11pt]{article}

\usepackage[margin=1in]{geometry}
\usepackage[T1]{fontenc}
\usepackage[utf8]{inputenc}
\usepackage{lmodern}
\usepackage{amsmath,amssymb}
\usepackage{booktabs}
\usepackage{array}
\usepackage{enumitem}
\usepackage{longtable}
\usepackage{xcolor}
\usepackage{hyperref}
\usepackage{listings}

\hypersetup{
    colorlinks=true,
    linkcolor=blue,
    urlcolor=blue
}

\lstset{
    basicstyle=\ttfamily\small,
    breaklines=true,
    columns=fullflexible,
    keepspaces=true,
    frame=single
}

\title{Composite Grid Representation}
\author{}
\date{}

\begin{document}
\maketitle

\section{Why the Project Uses a Composite Grid}

A normal grid MAPF implementation often stores only vertices and infers movement from adjacency. This project instead stores both:
\begin{itemize}[leftmargin=2em]
    \item \textbf{vertex slots}, and
    \item \textbf{transition slots}.
\end{itemize}

That choice matters because the project wants movement structure itself to be editable. The distinction between classical and cyclic mapping is not just a hidden rule in the solver. It is encoded directly in the map.

\section{Dimension Formula}

If the underlying base grid has dimensions
\[
R \times C,
\]
then the composite grid has dimensions
\[
(2R-1)\times(2C-1).
\]

\section{Slot Meaning}

The slot type is determined entirely by parity:
\begin{itemize}[leftmargin=2em]
    \item even row, even column: a real vertex,
    \item even row, odd column: a horizontal transition slot,
    \item odd row, even column: a vertical transition slot,
    \item odd row, odd column: a placeholder corner slot.
\end{itemize}

So a base-grid cell \((r,c)\) corresponds to composite-grid vertex position
\[
(2r,\,2c).
\]

\section{Stored Element Types}

The core enums live in \texttt{dev/core/composite\_elements.py}. The important conceptual categories are:
\begin{itemize}[leftmargin=2em]
    \item vertex values: free space or obstacle,
    \item horizontal transition values: left, right, bidirectional, or no transition,
    \item vertical transition values: up, down, bidirectional, or no transition,
    \item placeholder corner slots.
\end{itemize}

\section{Practical Consequences}

\subsection{Assignments and Semantic Masks}

When the code loads a binary obstacle matrix or a semantic campus mask, those values initially live in ordinary image/grid coordinates. Whenever the project needs to restrict spawning or identify zone membership inside the composite map, it converts raster coordinates to composite coordinates by doubling row and column indices.

\subsection{Mapping Becomes Explicit Data}

Because transitions are first-class cells, a cyclic map is not just a free-space grid with an orientation rule applied on the fly. It is a concrete composite grid in which each transition slot stores the allowed direction.

\subsection{Visualization Is Also Composite-Grid-Aware}

The renderer does not merely draw occupancy. It draws:
\begin{itemize}[leftmargin=2em]
    \item vertex cells,
    \item transition symbols,
    \item obstacle-adjacent side fills,
    \item placeholder-corner fills.
\end{itemize}
This is why the visualization logic is tightly coupled to the composite representation.

\section{Mental Model for Reading the Code}

When reading the repository, it is helpful to think of a composite map as a \emph{movement diagram} rather than a plain obstacle raster. The obstacles define which vertices are blocked, but the transition cells define what kinds of motion remain available between the free vertices.

\end{document}
